\textbf{Phần 1. Trắc nghiệm lựa chọn}
\setcounterpageref{ex}{0}

\begin{ex} %[5D10V59]
Một người đi xe máy khởi hành lúc $7\text{ giờ } 15\text{ phút}$ và đến nơi lúc $9\text{ giờ } 45\text{ phút}$. Biết quãng đường người đó đã đi dài $105\text{ km}$. Vận tốc của người đi xe máy là
\choice
{$40\text{ km/h}$}
{\True $42\text{ km/h}$}
{$45\text{ km/h}$}
{$50\text{ km/h}$}
\loigiai{
Thời gian người đó đi trên đường là:
$$9\text{ giờ } 45\text{ phút} - 7\text{ giờ } 15\text{ phút} = 2\text{ giờ } 30\text{ phút} = 2{,}5\text{ (giờ)}.$$
Vận tốc của người đi xe máy là:
$$v = 105 : 2{,}5 = 42\text{ (km/h)}.$$
\textbf{Đáp án của câu hỏi là B}
}
\end{ex}

\begin{ex} %[5D10V59]
Một vận động viên điền kinh chạy cự li $1\,500\text{ m}$ hết $4$ phút. Vận tốc chạy của vận động viên đó tính theo đơn vị mét trên giây là
\choice
{$6\text{ m/s}$}
{\True $6{,}25\text{ m/s}$}
{$6{,}5\text{ m/s}$}
{$5{,}75\text{ m/s}$}
\loigiai{
Đổi: $4\text{ phút} = 240\text{ giây}$.\\
Vận tốc chạy của vận động viên là:
$$v = 1\,500 : 240 = 6{,}25\text{ (m/s)}.$$
\textbf{Đáp án của câu hỏi là B}
}
\end{ex}

\begin{ex} %[5D10V59]
Một ô tô đi được quãng đường $162\text{ km}$ trong thời gian $3$ giờ $36$ phút. Vận tốc của ô tô đó là
\choice
{$42\text{ km/h}$}
{\True $45\text{ km/h}$}
{$48\text{ km/h}$}
{$50\text{ km/h}$}
\loigiai{
Đổi: $3\text{ giờ } 36\text{ phút} = 3 + \dfrac{36}{60} = 3{,}6\text{ (giờ)}$.\\
Vận tốc của ô tô là:
$$v = 162 : 3{,}6 = 45\text{ (km/h)}.$$
\textbf{Đáp án của câu hỏi là B}
}
\end{ex}

\begin{ex} %[5D10V59]
Biết một con chim cắt bay với vận tốc $90\text{ km/h}$. Vận tốc của con chim cắt đó đổi ra mét trên giây là
\choice
{$20\text{ m/s}$}
{\True $25\text{ m/s}$}
{$30\text{ m/s}$}
{$15\text{ m/s}$}
\loigiai{
Đổi: $90\text{ km} = 90\,000\text{ m}$; $1\text{ giờ} = 3\,600\text{ giây}$.\\
Vận tốc tính theo mét trên giây là:
$$90\,000 : 3\,600 = 25\text{ (m/s)}.$$
\textbf{Đáp án của câu hỏi là B}
}
\end{ex}

\begin{ex} %[5D10V59]
Một người đi bộ từ A đến B hết $45$ phút với vận tốc $4\text{ km/h}$. Sau đó người đó đi xe đạp từ B đến C dài $9\text{ km}$ hết $36$ phút. Vận tốc trung bình của người đó trên cả quãng đường từ A đến C là
\choice
{$8\text{ km/h}$}
{\True $8{,}89\text{ km/h}$}
{$9\text{ km/h}$}
{$8{,}5\text{ km/h}$}
\loigiai{
Đổi: $45\text{ phút} = 0{,}75\text{ giờ}$.\\
Quãng đường AB dài là:
$$4 \times 0{,}75 = 3\text{ (km)}.$$
Tổng quãng đường AC là:
$$3 + 9 = 12\text{ (km)}.$$
Tổng thời gian đi là:
$$45\text{ phút} + 36\text{ phút} = 81\text{ phút} = \dfrac{81}{60}\text{ giờ} = 1{,}35\text{ giờ}.$$
Vận tốc trung bình trên cả đoạn đường là:
$$12 : 1{,}35 \approx 8{,}89\text{ (km/h)}.$$
\textbf{Đáp án của câu hỏi là B}
}
\end{ex}

\begin{ex} %[5D10V59]
Một người đi xe máy từ nhà lúc $8\text{ giờ } 10\text{ phút}$, đi được $48\text{ km}$ thì tới cơ quan lúc $9\text{ giờ } 25\text{ phút}$. Dọc đường người đó phải dừng chờ đèn đỏ và đổ xăng mất $15\text{ phút}$. Vận tốc thực tế của xe máy khi chạy trên đường là
\choice
{$40\text{ km/h}$}
{\True $48\text{ km/h}$}
{$45\text{ km/h}$}
{$36\text{ km/h}$}
\loigiai{
Tổng thời gian từ lúc đi đến lúc tới nơi là:
$$9\text{ giờ } 25\text{ phút} - 8\text{ giờ } 10\text{ phút} = 1\text{ giờ } 15\text{ phút}.$$
Thời gian thực tế xe chạy là:
$$1\text{ giờ } 15\text{ phút} - 15\text{ phút} = 1\text{ giờ}.$$
Vận tốc thực tế của xe máy là:
$$v = 48 : 1 = 48\text{ (km/h)}.$$
\textbf{Đáp án của câu hỏi là B}
}
\end{ex}

\begin{ex} %[5D10V59]
Một đoàn tàu hỏa chạy qua một cây cầu dài $450\text{ m}$ hết $45$ giây. Biết đoàn tàu dài $150\text{ m}$. Vận tốc của đoàn tàu hỏa đó là
\choice
{$40\text{ km/h}$}
{\True $48\text{ km/h}$}
{$36\text{ km/h}$}
{$54\text{ km/h}$}
\loigiai{
Để vượt qua hoàn toàn cây cầu, đoàn tàu phải đi quãng đường bằng chiều dài cầu cộng chiều dài tàu:
$$s = 450 + 150 = 600\text{ (m)}.$$
Vận tốc của đoàn tàu theo $\text{m/s}$ là:
$$600 : 45 = \dfrac{40}{3}\text{ (m/s)}.$$
Đổi sang $\text{km/h}$:
$$\dfrac{40}{3} \times 3{,}6 = 48\text{ (km/h)}.$$
\textbf{Đáp án của câu hỏi là B}
}
\end{ex}

\begin{ex} %[5D10V59]
Một cano chạy xuôi dòng một khúc sông dài $45\text{ km}$ hết $1$ giờ $30$ phút. Biết vận tốc dòng nước là $3\text{ km/h}$. Vận tốc thực của cano khi nước yên lặng là
\choice
{$24\text{ km/h}$}
{\True $27\text{ km/h}$}
{$30\text{ km/h}$}
{$33\text{ km/h}$}
\loigiai{
Đổi: $1\text{ giờ } 30\text{ phút} = 1{,}5\text{ giờ}$.\\
Vận tốc ca nô khi xuôi dòng là:
$$v_{\text{xuôi}} = 45 : 1{,}5 = 30\text{ (km/h)}.$$
Vận tốc thực của ca nô là:
$$v_{\text{thực}} = v_{\text{xuôi}} - v_{\text{nước}} = 30 - 3 = 27\text{ (km/h)}.$$
\textbf{Đáp án của câu hỏi là B}
}
\end{ex}

\begin{ex} %[5D10V59]
Bạn Nam chạy quanh một sân vận động hình tròn có bán kính $50\text{ m}$ hết đúng $3$ phút $8$ giây. Lấy $\pi = 3{,}14$, vận tốc chạy của Nam xấp xỉ là
\choice
{$1{,}5\text{ m/s}$}
{\True $1{,}67\text{ m/s}$}
{$1{,}8\text{ m/s}$}
{$2\text{ m/s}$}
\loigiai{
Chu vi của sân vận động là:
$$50 \times 2 \times 3{,}14 = 314\text{ (m)}.$$
Đổi: $3\text{ phút } 8\text{ giây} = 188\text{ giây}$.\\
Vận tốc chạy của Nam là:
$$v = 314 : 188 \approx 1{,}67\text{ (m/s)}.$$
\textbf{Đáp án của câu hỏi là B}
}
\end{ex}

\begin{ex} %[5D10V59]
Một xe máy và một xe ô tô cùng khởi hành từ A đi về B. Quãng đường AB dài $120\text{ km}$. Xe ô tô đến B trước xe máy $30$ phút. Biết xe máy đi hết $3$ giờ. Vận tốc của ô tô lớn hơn vận tốc của xe máy là bao nhiêu ki-lô-mét trên giờ?
\choice
{$6\text{ km/h}$}
{\True $8\text{ km/h}$}
{$10\text{ km/h}$}
{$12\text{ km/h}$}
\loigiai{
Vận tốc của xe máy là:
$$120 : 3 = 40\text{ (km/h)}.$$
Thời gian ô tô đi từ A đến B là:
$$3\text{ giờ} - 30\text{ phút} = 2\text{ giờ } 30\text{ phút} = 2{,}5\text{ (giờ)}.$$
Vận tốc của ô tô là:
$$120 : 2{,}5 = 48\text{ (km/h)}.$$
Vận tốc ô tô lớn hơn vận tốc xe máy là:
$$48 - 40 = 8\text{ (km/h)}.$$
\textbf{Đáp án của câu hỏi là B}
}
\end{ex}

\begin{ex} %[5D10V59]
Trên một đoạn đường dốc dài $6\text{ km}$, một người đi xe đạp lên dốc hết $30$ phút và đi xuống dốc hết $15$ phút. Vận tốc trung bình của người đó trên cả chặng đường vừa lên dốc vừa xuống dốc là
\choice
{$15\text{ km/h}$}
{\True $16\text{ km/h}$}
{$18\text{ km/h}$}
{$20\text{ km/h}$}
\loigiai{
Tổng quãng đường cả lên và xuống dốc là:
$$6 + 6 = 12\text{ (km)}.$$
Tổng thời gian cả đi và về là:
$$30\text{ phút} + 15\text{ phút} = 45\text{ phút} = 0{,}75\text{ (giờ)}.$$
Vận tốc trung bình là:
$$v_{tb} = 12 : 0{,}75 = 16\text{ (km/h)}.$$
\textbf{Đáp án của câu hỏi là B}
}
\end{ex}

\begin{ex} %[5D10V59]
Vận tốc của một máy bay phản lực là $900\text{ km/h}$. Trong $1$ phút máy bay đó bay được quãng đường là
\choice
{$12\text{ km}$}
{\True $15\text{ km}$}
{$18\text{ km}$}
{$20\text{ km}$}
\loigiai{
Một giờ có $60$ phút.\\
Quãng đường máy bay bay được trong $1$ phút là:
$$900 : 60 = 15\text{ (km)}.$$
\textbf{Đáp án của câu hỏi là B}
}
\end{ex}

\textbf{Phần 2. Trắc nghiệm đúng sai}
\setcounterpageref{ex}{0}

\begin{ex} %[5D10V59]
Một người đi xe máy từ tỉnh A đến tỉnh B cách nhau $90\text{ km}$ hết $2$ giờ $15$ phút. Xét tính đúng/sai của các nhận định sau:
\choiceTF
{\True Thời gian người đó đi đổi ra giờ là $2{,}25$ giờ}
{\True Vận tốc của người đi xe máy là $40\text{ km/h}$}
{Nếu người đó tăng vận tốc thêm $5\text{ km/h}$ thì thời gian đi từ A đến B chỉ hết $1$ giờ $45$ phút}
{\True Đổi vận tốc của người đi xe máy ra mét trên phút được khoảng $667\text{ m/phút}$}
\loigiai{
\begin{itemchoice}
\item Đúng. $2\text{ giờ } 15\text{ phút} = 2{,}25\text{ giờ}$.
\item Đúng. $v = 90 : 2{,}25 = 40\text{ (km/h)}$.
\item Sai. Vận tốc mới là $40 + 5 = 45\text{ km/h}$. Thời gian đi là $90 : 45 = 2\text{ giờ}$ (không phải $1$ giờ $45$ phút).
\item Đúng. $40\text{ km/h} = 40\,000\text{ m} : 60\text{ phút} \approx 666{,}67\text{ m/phút} \approx 667\text{ m/phút}$.
\end{itemchoice}
}
\end{ex}

\begin{ex} %[5D10V59]
Hai bạn Hùng và Dũng tham gia chạy trên đường chạy thẳng dài $200\text{ m}$. Hùng chạy với vận tốc $5\text{ m/s}$, Dũng chạy với vận tốc $18\text{ km/h}$. Xét tính đúng/sai của các nhận định sau:
\choiceTF
{\True Vận tốc của Dũng đổi ra mét trên giây là $5\text{ m/s}$}
{\True Vận tốc của hai bạn bằng nhau}
{Hùng về đích trước Dũng $4$ giây}
{\True Cả hai bạn cùng về đích sau $40$ giây}
\loigiai{
\begin{itemchoice}
\item Đúng. $18\text{ km/h} = 18\,000\text{ m} : 3\,600\text{ s} = 5\text{ m/s}$.
\item Đúng. Cả hai bạn đều chạy với vận tốc $5\text{ m/s}$.
\item Sai. Vì vận tốc bằng nhau nên hai bạn về đích cùng lúc.
\item Đúng. Thời gian chạy là $200 : 5 = 40\text{ (giây)}$.
\end{itemchoice}
}
\end{ex}

\textbf{Phần 3. Câu trả lời ngắn}
\setcounterpageref{ex}{0}

\begin{ex} %[5D10V59]
Một chiếc ô tô đi từ Hà Nội đến Hải Dương trên quãng đường dài $72\text{ km}$ hết $1$ giờ $36$ phút. Vận tốc của chiếc ô tô đó là bao nhiêu ki-lô-mét trên giờ?
\shortans{45}
\loigiai{
Đổi: $1\text{ giờ } 36\text{ phút} = 1{,}6\text{ giờ}$.\\
Vận tốc của ô tô là:
$$72 : 1{,}6 = 45\text{ (km/h)}.$$
}
\end{ex}

\begin{ex} %[5D10V59]
Một người đi xe đạp đi hết quãng đường $36\text{ km}$ trong thời gian $2$ giờ $24$ phút. Vận tốc của người đi xe đạp đó là bao nhiêu ki-lô-mét trên giờ?
\shortans{15}
\loigiai{
Đổi: $2\text{ giờ } 24\text{ phút} = 2{,}4\text{ giờ}$.\\
Vận tốc của người đi xe đạp là:
$$36 : 2{,}4 = 15\text{ (km/h)}.$$
}
\end{ex}

\begin{ex} %[5D10V59]
Một vận động viên bơi cự li $200\text{ m}$ hết $2$ phút $30$ giây. Vận tốc bơi của vận động viên đó là bao nhiêu mét trên phút?
\shortans{80}
\loigiai{
Đổi: $2\text{ phút } 30\text{ giây} = 2{,}5\text{ phút}$.\\
Vận tốc bơi là:
$$200 : 2{,}5 = 80\text{ (m/phút)}.$$
}
\end{ex}

\begin{ex} %[5D10V59]
Một đoàn tàu hỏa chạy từ ga A lúc $6\text{ giờ } 30\text{ phút}$ và đến ga B lúc $10\text{ giờ } 30\text{ phút}$. Dọc đường tàu dừng đón khách ở các ga nhỏ mất tổng cộng $45\text{ phút}$. Biết quãng đường từ A đến B dài $195\text{ km}$. Vận tốc của đoàn tàu khi đang chạy là bao nhiêu ki-lô-mét trên giờ?
\shortans{60}
\loigiai{
Tổng thời gian từ ga A đến ga B là:
$$10\text{ giờ } 30\text{ phút} - 6\text{ giờ } 30\text{ phút} = 4\text{ giờ}.$$
Thời gian thực tế tàu chạy là:
$$4\text{ giờ} - 45\text{ phút} = 3\text{ giờ } 15\text{ phút} = 3{,}25\text{ (giờ)}.$$
Vận tốc của đoàn tàu là:
$$195 : 3{,}25 = 60\text{ (km/h)}.$$
}
\end{ex}

\begin{ex} %[5D10V59]
Một ca nô đi xuôi dòng quãng sông dài $60\text{ km}$ hết $2$ giờ. Biết vận tốc của dòng nước chảy là $4\text{ km/h}$. Hỏi vận tốc của ca nô khi đi ngược dòng trên quãng sông đó là bao nhiêu ki-lô-mét trên giờ?
\shortans{22}
\loigiai{
Vận tốc ca nô khi xuôi dòng là:
$$60 : 2 = 30\text{ (km/h)}.$$
Vận tốc thực của ca nô là:
$$30 - 4 = 26\text{ (km/h)}.$$
Vận tốc ca nô khi ngược dòng là:
$$26 - 4 = 22\text{ (km/h)}.$$
}
\end{ex}

\begin{ex} %[5D10V59]
Một con ngựa phi nước đại trên quãng đường $9\text{ km}$ hết $12$ phút. Vận tốc của con ngựa đó là bao nhiêu ki-lô-mét trên giờ?
\shortans{45}
\loigiai{
Đổi: $12\text{ phút} = 0{,}2\text{ giờ}$.\\
Vận tốc của con ngựa là:
$$9 : 0{,}2 = 45\text{ (km/h)}.$$
}
\end{ex}
